% !TeX root = ../../paper.tex

\subsection{The Four Challenges}
\label{sec:apparatus-challenges}

Each challenge pairs a virtual horror object with an everyday physical proxy chosen to match its expected texture. The proxies cost a few euros each; none of them is instrumented.

\paragraph{Bowl of eyes.} A rusted bowl on a wooden table is filled with realistic looking eyeballs. The player has to dig through them to find the coin at the bottom of the bowl. The physical proxy is a real metal bowl filled with rubber balls whose size matches the virtual eyes exactly; the virtual bowl is scaled at runtime to the measured dimensions of the physical one, so the rim is where the hand expects it.

\paragraph{Bucket of maggots.} A wooden bucket teems with writhing maggots, again with a coin buried underneath. The proxy is a bucket of gummy worms. Their damp, slightly sticky surface and the way a handful gives under pressure come close to what the visuals promise, at least closer than any dry material we tried.

\paragraph{Bloody pipe.} A pipe assembly bursts when the player pulls out the coin and sprays blood at them, driven by a particle system. The haptic counterpart is deliberately low-tech, in the tradition of Wizard-of-Oz setups: a researcher standing outside the play area sprays a fine mist of water from a bottle through a plastic tube aligned with the virtual pipe, timed to the burst. The spray lasts around five seconds and then stops together with the particle effect.

\paragraph{Cobwebs behind the closet.} The last coin hides behind a closet, in a mass of spider webs the player has to reach into. The proxy is synthetic Halloween spider webbing stretched behind an upright table. The material clings to skin and stretches when the hand pulls back---a sensation that is hard to fake with anything else and unpleasant even when one knows exactly what it is. The coin itself is hanging on a thin string, so the player feels the sensation of ripping it free when they attempt to remove it.

The set covers a deliberate range of contact types: immersing the whole hand (eyes, worms), a sudden stimulus arriving on the body (spray), and a clinging texture the player has to pull free from (webs). It also varies who controls the moment of contact---the player in three challenges, the experimenter in one.
